\begin{table}[h]
\vspace{-15pt}
\centering
\scriptsize
\caption{Risk Management Table outlining the main potential risks and mitigation strategies.}\label{table:risks}
\begin{tabular}{|p{1.6cm}|p{6.7cm}|p{7.5cm}|}
\hline
\textbf{Risk} & \textbf{Impact} & \textbf{Risk Management} \\ \hline

Rapid shifts in AI & Sudden advancements in AI technologies, such as foundational models or entirely new directions, could render existing approaches obsolete. & Actively engage in research communities, attend conferences, and continuously monitor industry and academic trends to adapt strategies accordingly. \\ \hline

Withdrawal of personnel & Loss of key personnel could delay progress, reduce expertise in critical areas, and impact overall project momentum. & Document critical knowledge and processes to ensure continuity. Recruit additional experienced competence in the event of key personnel loss. \\ \hline

Poor communication & Results in misaligned expectations, scope changes, and potential project derailment due to lack of clarity and feedback. & Establish regular communication channels and clearly define stakeholder roles and expectations. Keep stakeholders informed of project progress and risks. \\ \hline

Data privacy & Loss of stakeholder trust, legal liabilities, and compliance failures due to unauthorized access to sensitive data. & Implement stringent encryption, differential privacy, and regular security audits; ensure strict adherence to GDPR and other relevant regulations. \\ \hline

Regulatory Changes & Non-compliance with evolving regulations could delay project deployment or require significant rework. & Engage legal experts, track regulatory developments, and design adaptable compliance frameworks. \\ \hline

Ethical AI Concerns & Negative public perception or project delays due to ethical issues, such as bias or misuse of AI. & Develop and apply ethical guidelines, conduct bias audits, and engage diverse stakeholder groups in project oversight. \\ \hline

\end{tabular}
\vspace{-10pt}
\end{table}
